% Generated draft: requires mathematical and source review.
\documentclass[11pt]{article}
\usepackage[margin=25mm]{geometry}
\usepackage{fontspec}
\setmainfont{Latin Modern Roman}
\usepackage{amsmath,amssymb,mathtools,unicode-math}
\setmathfont{Latin Modern Math}
\usepackage{enumitem,needspace,xurl,hyperref}
\setlength{\parindent}{0pt}
\setlength{\parskip}{5pt}
\setlength{\emergencystretch}{3em}
\newcommand{\N}{\mathbb{N}}
\newcommand{\Z}{\mathbb{Z}}
\newcommand{\Q}{\mathbb{Q}}
\newcommand{\R}{\mathbb{R}}
\newcommand{\C}{\mathbb{C}}
\newcommand{\F}{\mathbb{F}}
\newcommand{\A}{\mathbb{A}}
\newcommand{\PP}{\mathbb{P}}
\newcommand{\E}{\mathbb{E}}
\newcommand{\eps}{\varepsilon}
\newcommand{\dd}{\,\mathrm{d}}
\begin{document}
\section*{UnsolvedMath ID: 30000858}
\textbf{Source problem code:} catalogue entry
\subsection{Definitions and mathematical statement}
Fix $d\geq2$ and put independent Bernoulli-$p$ open/closed states on the nearest-neighbor edges of $\Z^d$. Let $C(0)$ be the open cluster of the origin and define
\[
\theta(p)=\mathbb P_p(|C(0)|=\infty),\qquad
\chi(p)=\mathbb E_p|C(0)|=\sum_{x\in\Z^d}\mathbb P_p(0\leftrightarrow x),
\]
allowing $\chi=\infty$. Retain the source's threshold equality
\[
p_c=\inf\{p\in[0,1]:\theta(p)>0\}
=\sup\{p\in[0,1]:\chi(p)<\infty\}.
\]

For each fixed dimension, determine critical behavior at $p_c$ and near-critical behavior as $p\uparrow p_c$ or $p\downarrow p_c$. The requested program includes existence and values of the associated critical exponents, and their universality in the intended comparison class. For example, standard logarithmic exponent conventions for the two supplied observables would be
\[
\beta=\lim_{p\downarrow p_c}\frac{\log\theta(p)}{\log(p-p_c)},
\qquad
\gamma=-\lim_{p\uparrow p_c}\frac{\log\chi(p)}{\log(p_c-p)},
\]
when these limits exist. These formulas specify examples of the exponent question; they do not assert existence or supply predicted values. An exact power asymptotic with a finite nonzero amplitude is stronger than such a logarithmic limit.

The extract does not list all desired observables or exponents, prescribe exact-asymptotic versus logarithmic definitions, or identify the lattices, bond/site models or finite-range changes over which universality is to be tested. Thus the full general critical-behavior request is retained without inventing a numerical conjecture or a universal comparison across unspecified models. Despite the record's high-dimensional title, its mathematical scope explicitly includes every $d\geq2$.
\subsection{Short English statement}
Determine critical and near-critical behavior for independent nearest-neighbor bond percolation in every dimension at least two, including existence, values and universality of critical exponents. The exact exponent list, asymptotic strength and universality class are not supplied.
\subsection{Sources}
\begin{itemize}
\item Supplied report: \emph{Arbeitsgemeinschaft: Percolation} (2008), Oberwolfach Reports, DOI 10.4171/owr/2007/48, pp. 2851--2892; proposed year 2007. \url{https://doi.org/10.4171/owr/2007/48}.
\item Local UnsolvedMath registry, permanent ID 30000858, catalogue entry; created 25 August 2026, updated 16 September 2026. Its corrected\_verified note reports comparison with the official report. The supplied 21 August 2026 triage labels the record open but explicitly calls it an incomplete general critical-behavior question for which no single complete resolution can be assigned. No additional external reference is supplied, and no independent status assessment is asserted here.
\end{itemize}
\end{document}
